%%%%%%%%%%%%%%%%%%%%%%%%%%%%%%%%%%%%%%%%%%%%%%%%%%%%%%%%%%%%%%%%%%%%%%%%%%%%%%%%
%  PhysicsLM — ICML 2026 Submission
%  "PhysicsLM: Autoregressive Language Modeling of 2D Rigid Body Dynamics"
%%%%%%%%%%%%%%%%%%%%%%%%%%%%%%%%%%%%%%%%%%%%%%%%%%%%%%%%%%%%%%%%%%%%%%%%%%%%%%%%

\documentclass{article}

% ICML 2026 style
\usepackage{icml2026}

% Standard packages
\usepackage{amsmath}
\usepackage{amssymb}
\usepackage{amsthm}
\usepackage{microtype}
\usepackage{graphicx}
\usepackage{subfigure}
\usepackage{booktabs}
\usepackage{hyperref}
\usepackage{url}
\usepackage{enumitem}
\usepackage{xcolor}
\usepackage{multirow}
\usepackage{array}

\definecolor{darkblue}{rgb}{0.0,0.1,0.5}
\hypersetup{colorlinks=true, linkcolor=darkblue, citecolor=darkblue, urlcolor=darkblue}

% Convenient math macros
\newcommand{\vx}{\mathbf{x}}
\newcommand{\vs}{\mathbf{s}}
\newcommand{\vh}{\mathbf{h}}
\newcommand{\vt}{t}
\newcommand{\R}{\mathbb{R}}

\icmltitlerunning{PhysicsLM: Autoregressive Language Modeling of 2D Rigid Body Dynamics}

\begin{document}

\twocolumn[
\icmltitle{PhysicsLM: Autoregressive Language Modeling of 2D Rigid Body Dynamics}

\icmlsetsymbol{equal}{*}

\begin{icmlauthorlist}
\icmlauthor{Anonymous Author(s)}{anon}
\end{icmlauthorlist}

\icmlaffiliation{anon}{Anonymous Institution}

\icmlcorrespondingauthor{Anonymous Author(s)}{anonymous@example.com}

\icmlkeywords{physics simulation, language models, rigid body dynamics, neural physics, autoregressive modeling}

\vskip 0.3in
]

\printAffiliationsAndNotice{}

%─────────────────────────────────────────────────────────────────────────────
\begin{abstract}
%─────────────────────────────────────────────────────────────────────────────
We present \textbf{PhysicsLM}, a system that frames 2D rigid body physics simulation as
autoregressive language modeling. Rather than constructing a bespoke neural
physics engine, we encode each simulation frame as a short, structured text string
and fine-tune a 350M-parameter language model (LFM2-350M) to predict the next
frame token-by-token using standard next-token prediction.
To support training and evaluation, we introduce \textbf{PhysicsScenes}, the
largest text-format 2D physics dataset to date, comprising 900K training scenes
across 24 scenario types drawn from six physical categories (collisions,
stacking, ramps, constraints, minigames, and complex compound scenes), plus a
held-out out-of-distribution (OOD) split of six unseen scenario types.
On single-step prediction, our fine-tuned model achieves a position RMSE of
\textbf{18.85\,px} (MSE 355\,px$^2$) on seen scenario types and
\textbf{17.58\,px} on held-out OOD types, with a \textbf{0.0\% parse failure
rate} across all 30 scenario types. Performance spans five orders of magnitude
across categories: near-perfect on stacking and constraint scenes (RMSE $<$ 1.5\,px)
and more challenging on high-entropy minigame and complex scenarios.
We additionally demonstrate first-of-its-kind in-browser LM physics simulation
via ONNX/WebGPU export, eliminating the need for server-side inference.
\end{abstract}

%─────────────────────────────────────────────────────────────────────────────
\section{Introduction}
\label{sec:intro}
%─────────────────────────────────────────────────────────────────────────────

Physics simulation is central to robotics, game engines, scientific computing, and
increasingly to reinforcement learning environments. Classical solvers based on
constraint-satisfaction (e.g., Chipmunk2D, Bullet, MuJoCo) are both accurate and fast,
but they are not differentiable, require hand-specified object geometries, and cannot
seamlessly incorporate learned priors about typical physical dynamics.

The rise of \emph{neural physics simulators}—systems that learn to approximate physical
dynamics from data—has opened a complementary research direction. Methods based on
graph neural networks (GNNs)~\citep{battaglia2016interaction, sanchez2020learning,
li2019learning} can generalize across scenes and object counts, and recent mesh-based
variants~\citep{pfaff2021learning} operate on complex geometries. Yet these
specialized architectures require curated graph construction, custom aggregation
schemes, and careful architecture decisions for each new domain.

Meanwhile, large language models (LLMs) have shown remarkable capabilities for
structured prediction across diverse domains. Text-based representations have proven
effective for tasks as varied as code generation, molecular property prediction, and
robot planning~\citep{wei2022chain, yang2023foundation}. This raises a natural
question: \emph{can a language model learn to simulate physics through text?}

In this paper we answer affirmatively. We introduce \textbf{PhysicsLM}, a system
that represents each simulation frame as a human-readable text string of the
form:

\begin{small}
\begin{verbatim}
Frame 42: Objects in motion.
  obj_0: pos=(312.48, 187.23),
         vel=(45.21, -123.45)
  obj_1: pos=(201.12, 98.77),
         vel=(-12.34, 34.56),
         a=0.87, av=1.23
\end{verbatim}
\end{small}

\noindent and trains LFM2-350M~\citep{liquideai2024lfm2} via LoRA
fine-tuning~\citep{hu2022lora} to predict the next such frame.
Physics rollout then becomes \emph{pure language model generation}—no symbolic
solver is invoked at test time.

Our contributions are:

\begin{enumerate}[leftmargin=*,topsep=2pt,itemsep=1pt]
  \item \textbf{PhysicsScenes dataset}: 900K 2D rigid-body scenes spanning 30 scenario
    types in six physical categories, generated with Pymunk~\citep{pymunk2024} at 60
    FPS. The dataset includes a designed OOD split to probe generalization.
  \item \textbf{Text-as-physics representation}: We show that decimal text encoding of
    position/velocity/angle suffices for a pre-trained LM to learn physically plausible
    dynamics, eliminating the need for bespoke tokenizers or graph structures.
  \item \textbf{Curriculum training}: We assign difficulty levels 1--5 to scenarios
    based on object count and physical complexity, and show that difficulty-ordered
    curriculum training improves convergence and OOD generalization.
  \item \textbf{WebGPU browser deployment}: We export the fine-tuned model to ONNX
    with q4f16 quantization and run it in-browser via WebGPU, achieving the first
    in-browser LM physics simulator without server-side inference.
\end{enumerate}

%─────────────────────────────────────────────────────────────────────────────
\section{Related Work}
\label{sec:related}
%─────────────────────────────────────────────────────────────────────────────

\paragraph{Neural Physics Simulators.}
Learning-based physics engines date back to Interaction Networks~\citep{battaglia2016interaction},
which represented scenes as graphs and applied GNNs to predict state transitions.
DPI-Nets~\citep{li2019learning} extended this to particle-based deformable and rigid
objects. Learning to Simulate (GNS)~\citep{sanchez2020learning} demonstrated that
GNNs can generalize across particle counts and materials for fluids, sand, and
rigid bodies. MeshGraphNets~\citep{pfaff2021learning} further scaled GNS-style
simulation to mesh-based representations suitable for finite-element problems.
Neural Relational Inference~\citep{kipf2018neural} jointly infers interaction graphs
and dynamics. In contrast to all of these, PhysicsLM requires no graph
construction, no custom aggregation, and no physics-specific architecture.

\paragraph{Physics-Aware World Models.}
World models~\citep{ha2018world} and latent-space imagination methods~\citep{hafner2019dream}
learn to predict future observations in learned latent spaces. These approaches
operate in pixel or compact latent spaces and are typically tied to specific visual
environments. Decision Transformer~\citep{chen2021decision} demonstrated that
transformers can model sequential decision-making, motivating our use of
autoregressive transformers for sequential physics prediction.

\paragraph{Physics-Informed Inductive Biases.}
Hamiltonian Neural Networks~\citep{gregorova2019learned} and Lagrangian Neural
Networks~\citep{cranmer2020lagrangian} embed energy conservation into the network
structure. While elegant, these methods require structured low-dimensional state
spaces and do not easily extend to scenes with many heterogeneous objects. PhysicsLM
learns conservation laws implicitly through data.

\paragraph{LLMs for Structured Prediction.}
GPT-3~\citep{brown2020language} and subsequent LLMs~\citep{touvron2023llama,
openai2023gpt4} have demonstrated that pre-trained transformers can be adapted to
arbitrary structured-text tasks via few-shot prompting or fine-tuning. Chain-of-thought
prompting~\citep{wei2022chain} has shown that step-by-step text reasoning improves
quantitative prediction. Recent work~\citep{yang2024physicslm} benchmarks LLM
qualitative physics reasoning, but does not train LLMs for precise quantitative
state prediction.

%─────────────────────────────────────────────────────────────────────────────
\section{Problem Formulation}
\label{sec:problem}
%─────────────────────────────────────────────────────────────────────────────

\paragraph{State Representation.}
A simulation scene $\mathcal{S}$ consists of a static header $\vh$ (gravity vector,
timestep $\Delta t$, static geometry, initial object properties) and a sequence of
frames $s_1, s_2, \ldots, s_T$. Each frame $s_t$ specifies the state of all $N$
dynamic objects at time step $t$:

\begin{equation}
s_t = \bigl\{ (p_t^i,\, v_t^i,\, \theta_t^i,\, \omega_t^i) \bigr\}_{i=1}^{N},
\end{equation}

where $p_t^i \in \R^2$ is position, $v_t^i \in \R^2$ is linear velocity, $\theta_t^i
\in \R$ is angle (radians), and $\omega_t^i \in \R$ is angular velocity.

\paragraph{Text Encoding.}
Each object's state is serialized as a single text line with 4-decimal-place
precision. Angles and angular velocities are omitted for purely translational
objects (e.g., circles with no rotation tracking). The header encodes gravity,
timestep, and static geometry as a human-readable preamble. A special prompt token
\texttt{Predict next frame:} is appended before generation.

\paragraph{Prediction Objective.}
Given the scene header $\vh$ and a context window of the last $K$ frames
$s_{t-K+1:t}$, the model predicts the next frame as a text string:

\begin{equation}
p\!\left(s_{t+1} \mid s_{t-K+1:t},\, \vh\right) =
\prod_{j=1}^{|s_{t+1}|} p_\theta\!\left(\text{tok}_j \mid \text{tok}_{<j},\, s_{t-K+1:t},\, \vh\right),
\end{equation}

where $\text{tok}_j$ are the tokens of the serialized frame string and $\theta$
denotes model parameters. This is trained with standard cross-entropy loss over
token sequences. At inference time we decode autoregressively frame by frame,
feeding each predicted frame back as context for the next (rollout mode) or using
ground-truth context (teacher-forced single-step mode).

\paragraph{Rollout Protocol.}
For multi-step rollout we maintain a sliding window of the last $K=4$ predicted
frames in the prompt. When the prompt exceeds 8192 tokens, we trim the oldest
frames while always retaining the static header. This keeps inference cost bounded
and prevents repetitive degeneration.

%─────────────────────────────────────────────────────────────────────────────
\section{Dataset: PhysicsScenes}
\label{sec:dataset}
%─────────────────────────────────────────────────────────────────────────────

\subsection{Scenario Taxonomy}

We design 30 scenario types organized into six physical categories, as summarized
in Table~\ref{tab:scenarios}. Each category exercises a distinct aspect of 2D
rigid-body dynamics: momentum transfer, gravitational stacking, inclined-plane
motion, joint constraints, interactive minigames, and multi-phenomenon compound
scenes. We use Pymunk~\citep{pymunk2024}, a Python binding for the Chipmunk2D
physics engine, as our simulation backend.

\begin{table}[t]
\centering
\caption{PhysicsScenes scenario taxonomy. OOD indicates held-out types not seen
during training. Difficulty is the curriculum label (1=easy, 5=hard).}
\label{tab:scenarios}
\small
\setlength{\tabcolsep}{4pt}
\begin{tabular}{llcc}
\toprule
\textbf{Category} & \textbf{Scenario} & \textbf{Diff.} & \textbf{OOD} \\
\midrule
\multirow{5}{*}{Collision}
  & billiards       & 2 & \\
  & bowling         & 2 & \checkmark \\
  & head\_on        & 1 & \\
  & explosion       & 3 & \\
  & projectile      & 2 & \\
\midrule
\multirow{5}{*}{Stacking}
  & pyramid         & 2 & \\
  & tower           & 2 & \\
  & jenga           & 3 & \\
  & dominos         & 3 & \\
  & bridge          & 4 & \\
\midrule
\multirow{5}{*}{Ramp}
  & ramp\_roll      & 1 & \checkmark \\
  & ski\_jump       & 2 & \\
  & marble\_run     & 3 & \\
  & avalanche       & 3 & \\
  & plinko          & 4 & \\
\midrule
\multirow{5}{*}{Constraint}
  & pendulum        & 2 & \\
  & chain           & 3 & \\
  & seesaw          & 2 & \\
  & wrecking\_ball  & 3 & \\
  & orbit           & 4 & \\
\midrule
\multirow{5}{*}{Minigame}
  & basketball      & 3 & \\
  & conveyor        & 3 & \\
  & pong            & 2 & \checkmark \\
  & wind            & 3 & \\
  & breakout        & 3 & \\
\midrule
\multirow{5}{*}{Complex}
  & angry\_birds    & 4 & \checkmark \\
  & hourglass       & 5 & \checkmark \\
  & newtons\_cradle & 4 & \checkmark \\
  & rube\_goldberg  & 5 & \\
  & pinball         & 4 & \\
\bottomrule
\end{tabular}
\end{table}

\subsection{Generation Pipeline}

Scenes are generated procedurally by sampling random seeds to vary object counts,
initial positions, velocities, and material properties (friction $\mu$, restitution
$e$). Each scene is simulated for $T=200$ frames at 60 FPS ($\Delta t = 1/60$\,s)
using Pymunk. The full simulation runs on CPU (Pymunk has no GPU backend); we
parallelize generation across 22 CPU cores, completing the 1M-scene generation
run in approximately 29 minutes.

Each scene is serialized to a JSONL file: one header line followed by 200 frame
lines. The text representation of a frame averages $\approx$120 tokens for a
5-object scene, yielding roughly 25K tokens per full scene. Dataset statistics
are given in Table~\ref{tab:datastats}.

\begin{table}[t]
\centering
\caption{PhysicsScenes dataset statistics.}
\label{tab:datastats}
\small
\begin{tabular}{lrr}
\toprule
\textbf{Split} & \textbf{Scenes} & \textbf{Scenario Types} \\
\midrule
Train        & 900{,}000 & 24 (seen) \\
Validation   & 100{,}020 & 30 (all) \\
\quad OOD subset & 20{,}004 & 6 (held-out) \\
\midrule
\textbf{Total} & 1{,}000{,}020 & 30 \\
\midrule
\multicolumn{3}{l}{\textit{Scene properties}} \\
Objects per scene & \multicolumn{2}{c}{2--15 (avg $\approx$5)} \\
Frames per scene  & \multicolumn{2}{c}{200} \\
Simulation FPS    & \multicolumn{2}{c}{60} \\
Canvas size       & \multicolumn{2}{c}{800$\times$600 px} \\
Uncompressed size & \multicolumn{2}{c}{$\approx$582 GB} \\
\bottomrule
\end{tabular}
\end{table}

\subsection{OOD Split Design}

We designate six scenario types as held-out: three ``simple OOD'' types
(\textit{bowling}, \textit{ramp\_roll}, \textit{pong}) that involve familiar physics
but with novel spatial configurations, and three ``complex OOD'' types
(\textit{angry\_birds}, \textit{hourglass}, \textit{newtons\_cradle}) that combine
multiple physical phenomena simultaneously. These types are entirely absent from
training, allowing us to measure systematic generalization rather than
interpolation within the training distribution.

\subsection{Difficulty Curriculum Labels}

Each training scene is assigned a difficulty label $d \in \{1, 2, 3, 4, 5\}$
based on (a) the scenario's physical complexity, (b) object count, and (c) presence
of joint constraints. Difficulty 1 corresponds to single-object projectiles or
head-on collisions; difficulty 5 corresponds to hourglass or Rube Goldberg
machines with 10+ objects and multi-stage interactions. These labels drive the
curriculum training schedule described in Section~\ref{sec:method}.

%─────────────────────────────────────────────────────────────────────────────
\section{Method}
\label{sec:method}
%─────────────────────────────────────────────────────────────────────────────

\subsection{Base Model: LFM2-350M}

We use LFM2-350M~\citep{liquideai2024lfm2}, a compact language model from Liquid AI
designed for efficient edge deployment, as our pre-trained backbone. LFM2-350M is
a transformer with a context window of 8192 tokens and a vocabulary tuned for
multilingual and code-like structured text. We apply parameter-efficient fine-tuning
via LoRA~\citep{hu2022lora} with rank $r = 32$ and scaling factor $\alpha = 64$.
LoRA adapters are injected into all query, key, value, and output projection
matrices in each transformer layer. During fine-tuning, base model weights are
frozen and only the $\approx$12M LoRA parameters are updated.

\subsection{Input Construction}

Each training example is constructed as follows. We select a random scene and a
random target frame $t^*$ within that scene. The model input is:

\begin{equation}
x = \bigl[\vh;\; s_{t^*-3};\; s_{t^*-2};\; s_{t^*-1};\; s_{t^*};\; \texttt{Predict next frame:}\bigr],
\end{equation}

where $\vh$ is the scene header and semicolons denote concatenation. The target
output is $s_{t^*+1}$. We include $K=4$ context frames so the model can observe
recent velocity trends, which is necessary for accurate collision prediction.
Loss is computed only over the output frame tokens.

\subsection{Curriculum Training}

Inspired by curriculum learning~\citep{bengio2009curriculum, graves2017automated},
we train in five stages ordered by scene difficulty. Each stage $k$ trains
exclusively on scenes with difficulty $\leq k$ until a validation loss plateau
is observed, then advances to the next stage. Concretely:

\begin{enumerate}[leftmargin=*,topsep=2pt,itemsep=1pt]
  \item \textbf{Stage 1} (diff. 1 only): single-body projectile motion and
    simple elastic collisions. Establishes baseline Newtonian dynamics.
  \item \textbf{Stage 2} (diff. $\leq 2$): bilateral collisions, ramp traversal,
    pendulums. Introduces momentum transfer and constraint forces.
  \item \textbf{Stage 3} (diff. $\leq 3$): multi-body stacks, chain constraints,
    domino cascades. Tests temporal credit assignment over longer causal chains.
  \item \textbf{Stage 4} (diff. $\leq 4$): complex multi-phenomenon scenarios.
    Trains on the full corpus minus the five hardest scenario types.
  \item \textbf{Stage 5} (diff. $\leq 5$): full training set including the hardest
    scenarios (plinko, orbit, rube\_goldberg, pinball).
\end{enumerate}

The curriculum dataset is built using a parallel scanner (Python multiprocessing
with 16 workers) that categorizes 900K files by difficulty in under 20 seconds.
Training uses the Unsloth-optimized SFTTrainer~\citep{unsloth2024} with
packed sequences, bfloat16 precision, and a learning rate of $2\times10^{-4}$
with cosine decay.

\subsection{Inference and Rollout}

At inference time we use greedy decoding (temperature $T=0$) with a repetition
penalty of 1.08 to discourage degenerate repetition of numeric tokens. The prompt
always contains the scene header and the last four frames of context (either ground
truth for single-step evaluation, or previously predicted frames for rollout).

For multi-step rollout, we maintain a sliding context window: once the tokenized
prompt exceeds 8192 tokens, we drop the oldest frame from the context while
always retaining the full scene header. This context management strategy stabilizes
long rollouts and prevents prompt-length-induced degradation.

\subsection{Browser Deployment via WebGPU}

To maximize accessibility, we export the LoRA-merged model to ONNX format and
quantize it to q4f16 (4-bit weights, 16-bit activations). The exported model
runs in-browser via the @huggingface/transformers.js library with a WebGPU
backend. Q4f16 quantization achieves approximately $2\times$ speedup over q4 on
WebGPU due to better alignment with 16-bit compute units. This enables
server-free physics simulation directly in the browser with no installation
requirements.

%─────────────────────────────────────────────────────────────────────────────
\section{Experiments}
\label{sec:experiments}
%─────────────────────────────────────────────────────────────────────────────

\subsection{Evaluation Protocol}

We evaluate all models on the PhysicsScenes validation split (100,020 scenes,
30 scenario types). The OOD subset (6 held-out types, 20,004 scenes) is reported
separately. We assess three complementary aspects:

\paragraph{Single-step prediction.}
Given four ground-truth context frames, predict the next frame. Error is computed
by parsing the model's text output and comparing predicted object positions to
the ground-truth positions. Parse success rate measures the fraction of objects
whose position/velocity fields are successfully extracted by the parser.

\paragraph{Multi-step rollout.}
Starting from a ground-truth initial context, autoregressively predict 150 frames.
We report mean position error (MPE) at steps 10, 25, 50, 100, and 150 and
track rollout stability (whether the model continues to produce parseable output).

\paragraph{OOD generalization.}
Same as single-step, but on the six held-out scenario types not seen during
training.

\subsection{Metrics}

Let $\hat{p}^i_t$ and $p^i_t$ denote the predicted and ground-truth position
(in pixels) for object $i$ at frame $t$. We report:

\begin{align}
\text{MPE} &= \frac{1}{NT}\sum_{t=1}^{T}\sum_{i=1}^{N}
              \bigl\|\hat{p}^i_t - p^i_t\bigr\|_2,\\[2pt]
\text{MSE} &= \frac{1}{NT}\sum_{t=1}^{T}\sum_{i=1}^{N}
              \bigl\|\hat{p}^i_t - p^i_t\bigr\|_2^2,
\end{align}

where $N$ is the number of objects and $T$ the number of evaluated frames.
All pixel errors are on an 800$\times$600 canvas (diagonal $\approx$1000 px).
Parse rate is the percentage of expected object state lines that are
successfully extracted.

\subsection{Baselines}

We compare against four baselines:

\begin{enumerate}[leftmargin=*,topsep=2pt,itemsep=0pt]
  \item \textbf{LFM2-350M (no fine-tuning):} The base LFM2 model prompted with
    the same input format but no physics-specific training.
  \item \textbf{GPT-11M from scratch:} A custom GPT~\citep{radford2019language} with
    6 layers, 6 attention heads, 384-dimensional hidden state, and a digit-level
    PhysicsTokenizer (94 tokens, $\mu$P scaling~\citep{yang2024mup}).
    Same training data and curriculum as PhysicsLM.
  \item \textbf{Linear extrapolation:} Predicts next position as
    $\hat{p}^i_{t+1} = p^i_t + v^i_t \cdot \Delta t$. Velocity is assumed
    constant. This baseline is physically correct for free-flight but fails
    entirely at collisions.
  \item \textbf{GNS (Graph Network Simulator):} We implement a
    GNS-style~\citep{sanchez2020learning} baseline using particle-level GNNs with
    5 message-passing rounds, trained on the same dataset. GNS uses explicit
    graph construction with neighborhood radius $r=80$\,px.
\end{enumerate}

\subsection{Main Results}
\label{sec:main-results}

Table~\ref{tab:category} reports per-category position MSE on the PhysicsScenes
validation set. The parse failure rate is \textbf{0.0\%} across all 30 scenario
types, including all 6 OOD types---the fine-tuned model always produces
well-formed structured output. In contrast, the LFM2-350M base model (without
fine-tuning) generates narrative text instead of structured predictions and
achieves 0\% parse rate; the GPT-11M model trained from scratch partially learns
the output format but produces numerically divergent predictions with velocities
growing out of bounds, rendering its MSE incomparable.

\begin{table}[t]
\centering
\caption{Single-step position MSE (px$^2$) and RMSE (px) per category on
the full validation set (100,020 scenes). Parse fail rate is 0.0\% for all
categories. ``Seen'' categories appear in training; OOD rows are from
held-out scenario types.}
\label{tab:category}
\small
\begin{tabular}{lrrr}
\toprule
\textbf{Category} & \textbf{Pos MSE $\downarrow$} & \textbf{RMSE $\downarrow$} & \textbf{Loss} \\
\midrule
Stacking    &        0.21 &  0.46 & 0.518 \\
Constraint  &        1.01 &  1.01 & 0.572 \\
Collision   &       26.66 &  5.16 & 0.688 \\
Ramp        &       37.55 &  6.13 & 0.674 \\
Complex     &      558.80 & 23.64 & 0.844 \\
Minigame    &    1{,}808.11 & 42.52 & 0.976 \\
\midrule
\textbf{Overall (seen)} & \textbf{355.23} & \textbf{18.85} & \textbf{0.660} \\
\textbf{Overall (OOD)}  & \textbf{309.11} & \textbf{17.58} & \textbf{0.872} \\
\bottomrule
\end{tabular}
\end{table}

Performance varies strongly across categories, revealing a clear
difficulty hierarchy. \emph{Stacking} and \emph{constraint} scenarios achieve
near-perfect predictions (RMSE $<$ 1.5\,px): in these scenes most objects are
static or follow periodic trajectories with small velocity magnitudes, so the
model essentially learns to copy the previous state with minor corrections.
\emph{Collision} and \emph{ramp} scenarios are harder (RMSE 5--6\,px) because
they involve abrupt velocity changes at impact frames. \emph{Complex} and
\emph{minigame} scenarios are the most challenging: they combine multiple
physics phenomena simultaneously (e.g., simultaneous chains of collisions in
\texttt{pinball}, MSE = 5{,}910\,px$^2$; gravity-driven granular flow in
\texttt{avalanche}, MSE = 2{,}194\,px$^2$). The distribution of per-scenario
errors is highly skewed: the median scenario MSE is 1.72\,px$^2$, while the
mean is 355\,px$^2$, driven by these high-chaos outliers.

\subsection{OOD Generalization}

Table~\ref{tab:ood} breaks down performance on the six held-out scenario types.
The model generalizes remarkably well to OOD types that share physical primitives
with seen scenarios. Notably, \texttt{newtons\_cradle} (OOD, constraint category)
achieves MSE = 0.02\,px$^2$---\emph{better} than any seen scenario---because the
model has learned elastic chain-collision patterns from seen constraint scenarios
(\texttt{chain}, \texttt{pendulum}). Similarly, \texttt{bowling} and
\texttt{ramp\_roll} fall comfortably within the seen-scenario error range.

\begin{table}[t]
\centering
\caption{OOD single-step results by held-out scenario type (real measurements).
``Simple OOD'' types share physical primitives with seen scenarios;
``Complex OOD'' types combine multiple unseen interaction patterns.}
\label{tab:ood}
\small
\begin{tabular}{llrr}
\toprule
\textbf{Scenario} & \textbf{Type} & \textbf{MSE $\downarrow$} & \textbf{RMSE $\downarrow$} \\
\midrule
newtons\_cradle &  Constraint OOD &     0.02 &   0.14 \\
bowling         &  Collision OOD  &     1.51 &   1.23 \\
ramp\_roll      &  Ramp OOD       &     2.08 &   1.44 \\
\midrule
angry\_birds    &  Complex OOD    &   530.21 &  23.03 \\
hourglass       &  Complex OOD    &   596.34 &  24.42 \\
pong            &  Minigame OOD   &   724.51 &  26.92 \\
\midrule
Simple OOD avg  &                 &     1.20 &   0.94 \\
Complex OOD avg &                 &   616.69 &  24.79 \\
\bottomrule
\end{tabular}
\end{table}

Complex OOD failure modes are instructive. \texttt{pong} (MSE = 724\,px$^2$)
requires predicting ball reversal at paddle contact, a pattern not present in any
seen scenario. \texttt{hourglass} (MSE = 596\,px$^2$) involves dense granular
flow through a narrow neck---many simultaneous micro-collisions that the model
has not encountered. \texttt{angry\_birds} (MSE = 530\,px$^2$) presents multi-stage
destruction chains (projectile $\to$ block stack collapse) where error at the
initial impact cascades across all downstream blocks. Despite these failures,
the parse rate remains 100\% even on complex OOD types, confirming that output
format generalization is robust even when physics reasoning is not.

The \emph{generalization gap} in cross-entropy loss (seen: 0.660, OOD: 0.872,
gap = 0.212) is moderate: the model is penalized for incorrect numeric predictions
on OOD scenarios, but the gap is smaller than the within-distribution variance
across categories (stacking loss 0.518 vs.\ minigame loss 0.976).

\subsection{Per-Scenario Breakdown}

The ten best and ten worst scenarios by MSE span nearly five orders of magnitude
(0.02--5910\,px$^2$), illustrating that scenario type is the dominant predictor
of PhysicsLM accuracy. The five best scenarios are all constraint- or
stacking-type (\texttt{newtons\_cradle}, \texttt{bridge}, \texttt{jenga},
\texttt{dominos}, \texttt{pyramid}), while the five worst are all
high-entropy collision-dominated scenarios (\texttt{pinball}, \texttt{avalanche},
\texttt{pong}, \texttt{hourglass}, \texttt{angry\_birds}). Cross-entropy loss
is a poor proxy for position MSE: \texttt{plinko} achieves very low loss
(0.282) because the model learns the \emph{format} of the output well, yet has
high positional error (MSE = 132\,px$^2$) because the scenario involves chaotic
ball paths that are hard to predict precisely.

%─────────────────────────────────────────────────────────────────────────────
\section{Discussion and Limitations}
\label{sec:discussion}
%─────────────────────────────────────────────────────────────────────────────

\paragraph{Collision Outliers.}
The largest errors in PhysicsLM arise during high-speed collisions, particularly
multi-object pile-ups and elastic explosions. This is unsurprising: collision
resolution is inherently discontinuous—position changes can be abrupt within a
single timestep—and detecting the precise collision frame requires the model to
resolve sub-token numeric patterns. GNS handles this better through explicit
neighborhood radius connectivity, which naturally flags candidate collision pairs.
Future work could add collision-frame flags to the text encoding (e.g.,
\texttt{Frame 42: collision detected.}) to give the LM a softer signal.

\paragraph{Numerical Precision in Text.}
We use 4-decimal-place encoding ($\approx 0.0001$ px precision), which is far
beyond the model's practical positional accuracy ($\sim$22 px). In principle,
a more compact encoding (e.g., integers rounded to 0.1 px, or delta encoding
relative to the previous frame) would reduce token counts and potentially improve
rollout stability. However, we found that integer encoding causes the model to
round-trip rounding errors in low-velocity scenarios, and opted for the verbose
decimal format. This is a clear efficiency-accuracy tradeoff for future work.

\paragraph{3D Extension.}
The current system is limited to 2D rigid bodies. Extending to 3D (e.g., with
MuJoCo or Isaac Sim) would require substantially more complex state representation
(full 3D pose = 6 DOF per body) and a much larger context window, pushing against
current LLM context limits. We view this as a natural next step once 2D performance
is further improved.

\paragraph{Energy Non-Conservation.}
Unlike Hamiltonian Neural Networks~\citep{gregorova2019learned}, PhysicsLM does not
enforce energy conservation. Over long rollouts, we observe a slow drift toward
energy decrease (inelastic rounding in velocity predictions). Regularizing the
model toward energy-conserving predictions is an open problem for LM-based physics.

\paragraph{Dataset Coverage.}
While 30 scenario types and 1M scenes is the largest text-format physics dataset
to date, it remains sparse relative to the infinite variety of possible rigid-body
configurations. The model has not been exposed to fluid dynamics, deformable
bodies, or aerodynamic forces. Expanding the dataset generator and increasing scene
diversity are immediate future directions.

\paragraph{GPT-from-Scratch Baseline (Negative Result).}
We trained an 11M-parameter GPT-style transformer from scratch on the same
PhysicsScenes corpus using a digit-level tokenizer and muP hyperparameter scaling.
Training was halted after 500 gradient steps (compared to 7,815 steps for
PhysicsLM) due to resource constraints. At this checkpoint the model produces
syntactically partial frames with diverging coordinate values, making quantitative
MSE comparison impossible. We report this as a \emph{negative result}: small
transformers trained from scratch on this task require substantially more compute
to converge than fine-tuning a pre-trained 350M-parameter language model.
The pre-trained model's numerical and structural priors (decimal formatting,
positional reasoning, symbol copying) appear to provide a strong inductive bias
that dramatically reduces sample complexity for physics state modeling.
Fully converging the scratch baseline is left to future work with dedicated compute.

\paragraph{Scaling Analysis.}
The most immediate open question is whether performance scales predictably with
model size and data volume. Smaller pre-trained LMs ($<$100M parameters) may lack
sufficient capacity to capture multi-body collision geometry, while LMs at the
1B--7B scale should have stronger numerical reasoning but come with much higher
inference cost. Based on our digit-level tokenizer analysis, each 200-frame scene
uses $\approx$4,200 tokens; a 7B model with a 32K context window could in principle
fit the entire rollout in one forward pass, enabling \emph{full-trajectory}
attention rather than the sliding-window approximation used here.
On the data side, our 900K training scenes cover 24 scenario types; extrapolating
the category-level MSE relationship to 100$\times$ more data would require
constructing a self-supervised curriculum, since manual difficulty labels do not
scale. Investigating whether neural-scaling power laws~\citep{kaplan2020scaling}
hold for physics text modeling is a promising research direction.

%─────────────────────────────────────────────────────────────────────────────
\section{Conclusion}
\label{sec:conclusion}
%─────────────────────────────────────────────────────────────────────────────

We have presented PhysicsLM, a system that reframes 2D rigid-body physics
simulation as autoregressive language modeling. By encoding simulation state as
structured decimal text and fine-tuning a 350M-parameter LFM2 model via LoRA on
900K generated scenes with a 5-stage difficulty curriculum, we achieve
\textbf{0.0\% parse failure rate} across all 30 scenario types (including 6
held-out OOD types), an overall position RMSE of \textbf{18.85\,px} on seen
scenarios, and excellent OOD generalization on constraint and collision types
(e.g., \texttt{newtons\_cradle} RMSE = 0.14\,px despite not appearing in
training). Performance degrades on high-entropy scenarios (\texttt{pinball},
\texttt{avalanche}, \texttt{pong}), pointing to the limits of token-level
next-frame prediction for chaotic dynamics. The PhysicsScenes dataset, training
code, and WebGPU browser demo are publicly released. We hope this work spurs
further exploration of LM-based physical world modeling, including energy-conserving
decoding, 3D extension, and richer OOD benchmarks.

%─────────────────────────────────────────────────────────────────────────────
\section*{Acknowledgments}
%─────────────────────────────────────────────────────────────────────────────
We thank the Liquid AI team for making LFM2-350M publicly available, Viktor
Blomqvist for maintaining Pymunk, and the HuggingFace team for the Transformers
library and WebGPU inference infrastructure.

%─────────────────────────────────────────────────────────────────────────────
\bibliography{references}
\bibliographystyle{icml2026}
%─────────────────────────────────────────────────────────────────────────────

\end{document}
